\chapter{Visual Field Test (Perimetry)}

\section*{What Is This Test?}
A visual-field test measures the area a person can see while looking at a central target, including peripheral vision. Automated perimetry presents lights at different locations and records whether the patient sees them. It is a functional eye test used to detect and monitor loss of visual field.

\section*{Alternative Names}
Visual Field Examination; Perimetry; Automated Perimetry; Humphrey Visual Field Test; Peripheral Vision Test; Field of Vision Test

\section*{Why It Is Ordered}
Glaucoma evaluation and monitoring; assessment of optic-nerve or retinal disease; evaluation of unexplained visual difficulty; and assessment of neurological conditions affecting the visual pathways. It may be repeated to determine whether a defect is stable or progressing.

\section*{How This Test Is Performed}
The patient covers one eye, rests the chin and forehead on the instrument, looks at a central target, and presses a button whenever a light is seen. The other eye is then tested. Results are displayed as a map and are interpreted with visual acuity, eye-pressure measurements, optic-nerve findings, and examination quality.

\section*{Preparation, Risks, and Limitations}
Bring glasses if instructed and remain attentive throughout the test. The test is painless and has no physical risk, but fatigue, poor understanding, eyelid position, uncorrected refractive error, or inconsistent responses can affect reliability. A single abnormal field usually needs confirmation and clinical correlation.

\section*{References}
\begin{enumerate}
  \item MedlinePlus. Visual Field Examination and Glaucoma Tests.
  \item American Academy of Ophthalmology. Visual field testing.
\end{enumerate}

\clearpage
\section*{Understanding Results and Follow-up}
The laboratory report should identify the specimen, method, units, reference interval or decision limit, and whether the result is positive, negative, abnormal, or inconclusive. Reference intervals vary with age, sex, pregnancy, specimen type, assay, and laboratory; a result outside a range is not automatically a diagnosis, and a result within a range does not always exclude disease. Discuss unexpected results with the ordering clinician, who can decide whether repeat or confirmatory testing, treatment, or referral is appropriate. Do not change medicines or delay urgent care based on an isolated result.


\section*{Regulatory Status and Authorization Date}
Regulatory status applies to the specific assay, instrument, manufacturer, and intended use rather than to the test name alone. No single approval date can be assigned to this test category. For a commercial assay, verify the current product in the relevant regulator's database: the U.S. Food and Drug Administration (FDA) clearance, approval, or authorization; India's Central Drugs Standard Control Organisation (CDSCO) licence or permission; the European Union In Vitro Diagnostic Regulation (EU IVDR) CE certificate issued through the applicable conformity-assessment route; or the United Kingdom Medicines and Healthcare products Regulatory Agency (MHRA) registration and UKCA/CE status. Laboratory-developed tests may not have an individual FDA, CDSCO, EU IVDR, or MHRA approval date.
\clearpage
